% Generated draft: requires mathematical and source review.
\documentclass[11pt]{article}
\usepackage[margin=25mm]{geometry}
\usepackage{fontspec}
\setmainfont{Latin Modern Roman}
\usepackage{amsmath,amssymb,mathtools,unicode-math}
\setmathfont{Latin Modern Math}
\usepackage{enumitem,needspace,xurl,hyperref}
\setlength{\parindent}{0pt}
\setlength{\parskip}{5pt}
\setlength{\emergencystretch}{3em}
\newcommand{\N}{\mathbb{N}}
\newcommand{\Z}{\mathbb{Z}}
\newcommand{\Q}{\mathbb{Q}}
\newcommand{\R}{\mathbb{R}}
\newcommand{\C}{\mathbb{C}}
\newcommand{\F}{\mathbb{F}}
\newcommand{\A}{\mathbb{A}}
\newcommand{\PP}{\mathbb{P}}
\newcommand{\E}{\mathbb{E}}
\newcommand{\eps}{\varepsilon}
\newcommand{\dd}{\,\mathrm{d}}
\begin{document}
\section*{UnsolvedMath ID: 30002984}
\textbf{Source problem code:} catalogue entry
\subsection{Definitions and mathematical statement}
\textbf{Scope and conventions.} The supplied record poses a research programme about hyperkähler manifolds of dimension greater than two; it does not specify enough hypotheses to determine a single classification problem or theorem.

In one common complex-geometric convention, an irreducible holomorphic symplectic manifold is a simply connected compact Kähler manifold \(X\) such that
\[
H^0(X,\Omega_X^2)=\C\sigma,
\]
where \(\sigma\) is an everywhere nondegenerate holomorphic two-form. Under this convention, \(\dim_{\C}X=2n\), and the dimension restriction becomes \(n\geq 2\). The source does not specify whether its term ``hyperkähler manifold'' means precisely this class, a broader class, or whether dimension is complex or real. This definition therefore explains a possible convention without adding it to the supplied hypotheses.

A common framework for a one-parameter degeneration is a proper flat family
\[
\pi:\mathcal X\longrightarrow\Delta,
\qquad \Delta^*=\Delta\setminus\{0\},
\]
over a complex disc, smooth over \(\Delta^*\), whose fibres over \(\Delta^*\) are manifolds in the class under consideration. The fibre \(\mathcal X_0\) is the limiting fibre. The source does not specify this framework, permitted singularities, projectivity, or the equivalence relation used to classify such families or limiting fibres.

For a specified object \(Y\) and deformation category, a versal deformation is a deformation over a pointed base germ from which every sufficiently small deformation of \(Y\) in that category is induced by base change; the inducing map need not be unique. A partial compactification of a moduli space \(M\) generally means a space \(\overline M\) containing \(M\) as a dense open subspace and adding boundary objects, without necessarily requiring \(\overline M\) to be proper. The intended deformation category, moduli problem, boundary objects, and additional properties of \(\overline M\) are unspecified here.

\textbf{Source question.} For hyperkähler manifolds of dimension greater than two, with the preceding conventions still to be fixed:
\begin{enumerate}
\item Classify their degenerations, after specifying the families or limiting objects to be classified and the equivalence relation.
\item Describe the known examples of such degenerations.
\item Describe what is known about their versal deformation spaces, after specifying which objects are being deformed.
\item Determine when their moduli spaces admit partial compactifications, after specifying the moduli problem and the required boundary and compactification properties.
\end{enumerate}
These four requests retain the source's scope. The missing specifications prevent their replacement by a uniquely determined quantified mathematical assertion.
\subsection{Short English statement}
The source asks for a classification and examples of degenerations of hyperkähler manifolds of dimension greater than two, information about versal deformation spaces, and conditions for partial compactifications of their moduli spaces. Its geometric conventions and moduli conditions are not fixed, so this remains a research programme requiring clarification and human mathematical review.
\subsection{Sources}
\begin{itemize}
\item Supplied local UnsolvedMath registry record, permanent ID \(30002984\), title: ``Degenerations and Compactifications of Hyperkähler Manifolds''; source problem code \texttt{catalogue entry}, in the Oberwolfach Reports Open Problems collection. No individual proposer is supplied. The proposed year is recorded as 2015.
\item Supplied workshop citation: \emph{Mini-Workshop: Singular Curves on \(K3\) Surfaces and Hyperkähler Manifolds} (2016), Oberwolfach Reports, pp. 2939--2967, DOI 10.4171/owr/2015/51. \url{https://doi.org/10.4171/owr/2015/51}. The supplied description says that these reports document workshop discussions and open problems. The citation's 2016 label differs from the recorded proposed year 2015; the relationship has not been independently checked.
\item Supplied literature reference: Kollar--Laza--Sacca--Voisin, \emph{Degenerations of Hyper-Kahler Manifolds}. \url{https://doi.org/10.5802/aif.3228}.
\item Supplied literature reference: Huybrechts--Mauri, \emph{Type II Degenerations}. \url{https://arxiv.org/abs/2108.01587}.
\item Supplied literature reference: \emph{Degenerations of Generalized Kummer Varieties}. \url{https://arxiv.org/abs/2604.14890}. No authors are supplied for this reference.
\item Literature assessment dated exactly \texttt{2026-08-21}: the supplied curation labels the problem ``partially solved'' and reports that Kulikov-type models, monodromy constraints, and smoothing results provide substantial control, while classification of higher-dimensional degenerations, all examples, and versal spaces remains far from complete. This is supplied dated curation derived from the OWR record, not an independently verified assessment; time-sensitive status and cited claims require independent verification.
\item The supplied registry labels the statement \texttt{corrected\_verified} and reports that it was checked against the official Oberwolfach report, preserving its hypotheses, notation, and open target. This records the supplied verification claim; no source was visited or independently verified in preparing this draft.
\end{itemize}
\end{document}
